\section{Giải thích cho lời giải thích}
\label{sec:ch11-giai-thich-cho-loi-giai-thich}

\index{attribution!nối tầng|(}%
Tầng thứ nhất là thứ Phần~II đã dựng xong. Một phương pháp attribution nhận một
mô hình $f$ và một đầu vào $x$, rồi trả về cho mỗi \tn{đặc trưng}{feature} $i$ một số
$\varphi_i(x,f)$ nói đặc trưng ấy đóng góp bao nhiêu vào đầu ra, theo nghĩa mà
định nghĩa~\ref{def:ch06-attribution} đặt. LIME, SHAP, Integrated Gradients và
Grad-CAM đều trả về một bảng số như vậy, mỗi đặc trưng một số.

Một bảng số như vậy không nói được hai đặc trưng ảnh hưởng lẫn nhau ra sao. Nếu
mô hình chỉ dùng đặc trưng thứ hai khi đặc trưng thứ nhất có mặt, thì quan hệ ấy
phải nằm đâu đó trong hai con số, nhưng hai con số không có chỗ để chứa nó: mỗi
số gắn với một đặc trưng, và quan hệ thì gắn với một cặp. Bài báo gọi những
tương tác ấy bằng hai tên đối nhau, synergy và antisynergy, và chỉ giải nghĩa
tên thứ hai: antisynergy là trường hợp attribution riêng của hai đặc trưng
triệt nhau khi cộng lại, tức đúng chỗ mà một phương pháp bậc nhất không nhìn
thấy gì~\autocite{p23metagame}.

Lối ra mà lĩnh vực chọn cho vấn đề ấy là lối mà chương này đọc, và nó có tên
riêng trong bài. Một \tn{attribution nối tầng}{serial attribution} lấy chính
phương pháp attribution áp lần thứ hai, lần này lên đầu ra của lần thứ nhất.
Cụ thể, hàm được giải thích ở lần thứ hai không còn là mô hình $f$ nữa mà là
hàm $\varphi_i(\cdot,f)$, tức hàm gán cho mỗi đầu vào con số mà lần thứ nhất đã
gán cho đặc trưng $i$:
%
\begin{equation}
  \label{eq:ch11-noi-tang}
  \varphi_j\bigl(x,\ \varphi_i(\cdot,f)\bigr).
\end{equation}
%
Đọc phương trình~\ref{eq:ch11-noi-tang} từ trong ra ngoài: bên trong là
attribution của đặc trưng $i$ cho mô hình, bên ngoài là attribution của đặc
trưng $j$ cho cái attribution ấy. Kết quả là một con số cho mỗi cặp $(i,j)$,
tức đúng chỗ chứa mà bảng số một tầng không có.

Ý tưởng ấy không mới trong bài báo này. Bài dẫn nó về Janizek, Sturmfels và Lee
năm 2021, những người đặt cho nó cái tên mà chương này mượn làm tên mục, giải
thích cho lời giải thích, và xây dựng phiên bản dùng Integrated Gradients hai
lần dưới tên integrated Hessians; phiên bản dùng \tn{giá trị Shapley}{Shapley value} hai lần thì
mang tên serial Shapley value và thuộc về Lundstrom cùng
Razaviyayn~\autocite{p23metagame}. Sách không đọc hai bài ấy và thuật lại chúng
qua lời của bài báo, theo đúng cách chương~\ref{ch:chi-so-khong-do-do-trung-thuc}
đã làm với những công trình mà bài báo neo nhắc tới.

Ký hiệu của bài lệch với ký hiệu của sách ở hai chỗ, và
phụ lục~\ref{app:ky-hieu} ghi phép đối chiếu cho người đọc mở bài ra.

Cấu trúc ấy trả lời được câu hỏi về cặp đặc trưng, nhưng nó mở ra một câu hỏi
mà mục sau dành cả mục để trả lời. Nếu con số ở tầng hai được sinh ra bằng cùng
một thủ tục đã sinh ra con số ở tầng một, thì nó cũng thừa hưởng mọi tính chất
của thủ tục ấy, kể cả những tính chất không mong muốn. Bài báo chỉ ra rằng có
một tính chất không mong muốn như thế, và chỉ ra nó trên một hàm mà người ta
biết trước đáp án đúng.
\index{attribution!nối tầng|)}%
